\chapter{Contexte général et état de l'art}
\label{chap:contexte}

\callout{En bref}{Ce chapitre situe le projet dans son cadre académique, rappelle
les fondements de la Business Intelligence et de sa chaîne décisionnelle, expose
les limites des outils existants, puis dresse l'état de l'art des briques
technologiques mobilisées : grands modèles de langage, agents d'IA et protocole MCP.}

\section{Cadre du projet}
Ce travail s'inscrit dans le cadre du \textbf{Projet de Fin d'Année (PFA)} de
la filière \textit{\mafiliere} à l'\textbf{\monetablissement}. Le PFA constitue
un jalon pédagogique majeur : il vise à confronter les étudiants à un problème
d'ingingénierie réaliste, à mobiliser de façon intégrée les connaissances
acquises (programmation, science des données, génie logiciel, intelligence
artificielle) et à développer des compétences de conception, de gestion de
projet et de communication scientifique.

Le sujet retenu se situe à l'intersection de trois domaines en pleine
expansion : la \textbf{Business Intelligence}, l'\textbf{ingénierie des données}
et l'\textbf{IA générative}. Il répond à une préoccupation concrète des
organisations : réduire le temps et l'expertise nécessaires pour transformer
des données brutes en connaissances décisionnelles.

\section{La Business Intelligence et la chaîne décisionnelle}

\subsection{Définition et finalité}
La \textbf{Business Intelligence} (informatique décisionnelle) désigne
l'ensemble des méthodes, des processus et des technologies permettant de
collecter, d'organiser, d'analyser et de restituer les données d'une
organisation afin d'éclairer la prise de décision. Son objectif n'est pas la
donnée pour elle-même, mais la \emph{connaissance actionnable} : répondre aux
questions « que s'est-il passé ?\,», « pourquoi ?\,» et « que faire ?\,».

\subsection{La chaîne décisionnelle}
La création de valeur à partir des données suit classiquement une chaîne en
plusieurs maillons, illustrée par la figure~\ref{fig:chaine}.

\begin{figure}[ht]
\centering
\begin{tikzpicture}[
  node distance=4pt and 14pt,
  box/.style={rectangle,rounded corners=4pt,draw=accent,fill=accentsoft,
    text=ink,minimum height=1.3cm,text width=2.4cm,align=center,font=\small\bfseries},
  arr/.style={-{Stealth[length=3mm]},accentdark,thick}]
  \node[box] (a) {Données brutes};
  \node[box,right=of a] (b) {Ingestion \& validation};
  \node[box,right=of b] (c) {Préparation \& nettoyage};
  \node[box,right=of c] (d) {Analyse \& KPIs};
  \node[box,right=of d] (e) {Restitution \& décision};
  \draw[arr] (a)--(b); \draw[arr] (b)--(c); \draw[arr] (c)--(d); \draw[arr] (d)--(e);
\end{tikzpicture}
\caption{La chaîne décisionnelle classique de la Business Intelligence.}
\label{fig:chaine}
\end{figure}

Chaque maillon mobilise des savoir-faire distincts :
\begin{itemize}
  \item \textbf{Ingestion} : récupération des données depuis des sources
        hétérogènes (fichiers, bases de données, API) et contrôle de leur
        intégrité ;
  \item \textbf{Préparation} : nettoyage, traitement des valeurs manquantes,
        typage, normalisation ;
  \item \textbf{Analyse} : définition et calcul des indicateurs clés de
        performance (KPI), exploration statistique ;
  \item \textbf{Restitution} : conception de tableaux de bord et de rapports
        destinés aux décideurs.
\end{itemize}

\subsection{Notions d'indicateurs, dimensions et mesures}
La modélisation décisionnelle distingue traditionnellement les
\textbf{mesures} (grandeurs numériques agrégeables : chiffre d'affaires,
quantités, buts marqués) des \textbf{dimensions} (axes d'analyse catégoriels
ou temporels : région, produit, date) selon lesquels les mesures sont
ventilées. Un \textbf{KPI} (\textit{Key Performance Indicator}) est une
mesure synthétique, souvent agrégée selon une ou plusieurs dimensions,
choisie pour sa pertinence métier. Cette distinction, fondamentale, est au
cœur de l'approche adaptative développée dans ce projet
(chapitre~\ref{chap:realisation}).

\section{Problématique : les limites des outils BI classiques}
Malgré leur richesse fonctionnelle, les solutions de BI du marché présentent
plusieurs limites structurantes :

\begin{description}
  \item[Forte dépendance à l'expertise.] La conception d'un tableau de bord
    pertinent suppose une bonne compréhension métier \emph{et} technique.
    L'analyste reste un goulot d'étranglement.
  \item[Travail répétitif et chronophage.] Chaque nouveau jeu de données
    impose de reparcourir la chaîne décisionnelle : exploration, choix des
    indicateurs, construction des graphiques, rédaction.
  \item[Spécialisation par cas d'usage.] Les solutions sont souvent
    paramétrées pour un domaine précis (BI commerciale, BI RH), avec des
    hypothèses figées sur la nature des colonnes.
  \item[Restitution manuelle.] La rédaction du rapport narratif — pourtant
    essentielle à la décision — demeure entièrement humaine.
\end{description}

Ces limites motivent la recherche d'une plateforme \textbf{généraliste} et
\textbf{automatisée}, capable de s'adapter à n'importe quel schéma de données
et de produire, sans intervention humaine, l'ensemble des livrables
décisionnels. Le tableau~\ref{tab:comparaison} synthétise l'écart entre
l'approche traditionnelle et celle que nous proposons.

\begin{table}[ht]
\centering\small
\caption{Comparaison entre la BI traditionnelle et l'approche proposée.}
\label{tab:comparaison}
\begin{tabularx}{\linewidth}{@{}lYY@{}}
\toprule
\textbf{Critère} & \textbf{BI traditionnelle} & \textbf{Approche proposée} \\
\midrule
Choix des KPIs & Manuel, par un expert & Automatique et adaptatif \\
Conception des graphiques & Manuelle & Automatique selon les rôles \\
Rédaction du rapport & Humaine & Générée par agents IA \\
Adaptation au domaine & Paramétrage spécifique & Aucune hypothèse métier \\
Délai d'obtention & Heures à jours & Quelques secondes \\
Véracité des chiffres & Confiance humaine & Vérification automatique \\
\bottomrule
\end{tabularx}
\end{table}

\section{Objectifs et périmètre}
Le périmètre du projet couvre l'intégralité de la chaîne décisionnelle, de
l'ingestion d'un fichier tabulaire jusqu'à la production d'un rapport
décisionnel. Les choix structurants sont les suivants :
\begin{itemize}
  \item support des formats \texttt{CSV}, \texttt{Excel} et \texttt{JSON} ;
  \item \textbf{aucune} configuration spécifique au domaine métier ;
  \item génération automatique d'un tableau de bord interactif et d'un
        rapport (web + PDF) ;
  \item exécution locale, reproductible, avec un magasin d'artefacts par
        exécution.
\end{itemize}
Sont hors périmètre : la connexion temps réel à des bases de données de
production, la gestion multi-utilisateurs et le déploiement à l'échelle, qui
constituent des perspectives (voir la conclusion générale).

\section{État de l'art}

\subsection{Les grands modèles de langage (LLM)}
Un \textbf{grand modèle de langage} est un réseau de neurones, fondé sur
l'architecture \textit{Transformer}~\cite{vaswani2017attention}, entraîné sur
d'immenses corpus textuels pour prédire le mot suivant. Au-delà de la
génération de texte, ces modèles manifestent des capacités de raisonnement, de
synthèse et de \emph{suivi d'instructions}. Une avancée déterminante pour notre
projet est l'\textbf{appel d'outils} (\textit{tool use} / \textit{function
calling})~: le modèle peut décider d'invoquer une fonction externe avec des
arguments structurés, puis exploiter le résultat. C'est précisément ce
mécanisme qui permet à un LLM de piloter un système logiciel plutôt que de se
contenter de produire du texte.

Nous mobilisons des modèles performants tels que \textbf{Llama~3.3}~\cite{llama3}
servis par l'infrastructure \textbf{Groq} (inférence à très faible latence), ou
des modèles exécutés localement via \textbf{Ollama}. Les LLM présentent
néanmoins deux faiblesses notoires : la tendance à l'\textbf{hallucination}
(production d'informations erronées mais plausibles) et la difficulté à
manipuler fidèlement de grandes tables. Ces limites guident notre choix
architectural : confier les calculs à un moteur déterministe et réserver aux
LLM la curation et la rédaction.

\subsection{Les agents d'intelligence artificielle}
Un \textbf{agent} est un système qui perçoit un environnement, raisonne et
agit pour atteindre un objectif. Appliqué aux LLM, le paradigme agentique
consiste à doter un modèle d'une \textbf{boucle} : observer l'état, décider
d'une action (souvent un appel d'outil), exécuter, observer le résultat, et
recommencer jusqu'à l'accomplissement de la tâche. Les approches multi-agents
décomposent un problème complexe en \textbf{rôles spécialisés} qui coopèrent,
à l'image d'une équipe humaine. Cette spécialisation améliore la fiabilité,
la lisibilité et la maintenabilité, chaque agent disposant d'un contexte et
de permissions restreints à sa mission.

\subsection{Le protocole MCP (Model Context Protocol)}
Le \textbf{Model Context Protocol}~\cite{mcp} est un standard ouvert,
introduit par Anthropic, qui normalise la manière dont une application expose
ses \textbf{outils}, ses \textbf{ressources} et ses \textbf{données} à des
modèles de langage. À la manière d'un « port universel » entre les LLM et le
monde logiciel, MCP définit un format commun pour décrire un outil (nom,
description, schéma JSON des arguments) et pour l'invoquer.

Dans notre projet, MCP joue un double rôle structurant : il sert de
\textbf{registre central des capacités} du système (chargement de données,
profilage, calcul d'indicateurs, génération de graphiques, etc.) et de
\textbf{mécanisme de contrôle d'accès}. Chaque agent se voit attribuer un
sous-ensemble d'outils autorisés ; toute tentative d'appel hors de ce
périmètre est refusée par le serveur. Ce principe de \emph{moindre privilège}
renforce la sûreté et clarifie les responsabilités.

\subsection{Modélisation dimensionnelle et entrepôts de données}
La théorie de la \textbf{modélisation dimensionnelle}, formalisée notamment par
Kimball~\cite{kimball2013}, structure les données décisionnelles autour de
\emph{tables de faits} (contenant les mesures) et de \emph{tables de dimensions}
(contenant les axes d'analyse). Cette séparation conceptuelle entre
\textbf{faits/mesures} et \textbf{dimensions} irrigue directement notre moteur de
profilage : plutôt que d'imposer un schéma en étoile prédéfini, le système
\emph{infère} dynamiquement le rôle de chaque colonne. Cette inférence
automatique des rôles constitue la transposition « sans configuration » des
principes de la modélisation dimensionnelle à un contexte généraliste.

\subsection{Outils de visualisation et frameworks web}
La restitution s'appuie sur \textbf{Plotly}~\cite{plotly}, une bibliothèque de
visualisation interactive produisant des graphiques en HTML/JavaScript
autonomes, et sur \textbf{FastAPI}~\cite{fastapi}, un framework web Python
moderne et performant pour l'interface utilisateur. La manipulation des
données repose sur \textbf{pandas}~\cite{pandas}, référence de l'analyse de
données tabulaires en Python.

\section{Méthodologie de travail}
Le projet a été conduit selon une démarche \textbf{itérative et incrémentale}.
Une première itération a établi le socle (moteur déterministe, serveur MCP,
sept agents, interface) ; les itérations suivantes ont enrichi le système
(robustesse, vérification des chiffres, suivi des coûts, cartes
géographiques, détection d'anomalies, comparaison de segments) et amélioré
l'interface. Chaque incrément a été \textbf{validé par des tests}
automatisés et par des exécutions de bout en bout sur des jeux de données
réels, garantissant la non-régression. Cette approche a permis de livrer en
permanence un système fonctionnel et de prioriser les fonctionnalités à plus
forte valeur.

\section{Gestion et planification du projet}
Le projet a été planifié en \textbf{cinq phases} successives mais partiellement
recouvrantes, conformément à la démarche itérative. La
figure~\ref{fig:gantt} en présente le diagramme de Gantt simplifié.

\begin{figure}[ht]
\centering
\begin{tikzpicture}[x=0.95cm,y=0.62cm,font=\footnotesize]
  % axe temps (semaines)
  \foreach \x in {0,...,12} \draw[muted!40] (\x,0.4) -- (\x,-5.2);
  \foreach \x/\lab in {0/S1,2/S3,4/S5,6/S7,8/S9,10/S11,12/S13}
     \node[muted,font=\scriptsize] at (\x,0.8) {\lab};
  % barres : nom, debut, duree, ligne
  \def\bar#1#2#3#4{\fill[accent,rounded corners=2pt] (#2,-#4+0.2) rectangle (#2+#3,-#4-0.2);
     \node[anchor=east,text=ink] at (-0.2,-#4) {#1};}
  \bar{Étude \& état de l'art}{0}{2.5}{1}
  \bar{Conception}{2}{2.5}{2}
  \bar{Développement du socle}{3.5}{4}{3}
  \bar{Enrichissement \& fiabilité}{7}{4}{4}
  \bar{Tests \& rédaction}{9.5}{3}{5}
\end{tikzpicture}
\caption{Diagramme de Gantt simplifié du déroulement du projet.}
\label{fig:gantt}
\end{figure}

\begin{description}
  \item[Étude \& état de l'art.] Cadrage du sujet, revue des concepts (BI, LLM,
    agents, MCP) et choix technologiques.
  \item[Conception.] Définition des besoins, de l'architecture en couches et
    de la répartition des rôles entre agents.
  \item[Développement du socle.] Moteur déterministe, serveur MCP, sept agents,
    interface web et premier pipeline de bout en bout.
  \item[Enrichissement \& fiabilité.] Filets de sécurité, vérification des
    chiffres, suivi des coûts, cartes géographiques et analyses avancées.
  \item[Tests \& rédaction.] Consolidation des tests, étude de cas et rédaction
    du présent rapport.
\end{description}

\callout{Synthèse du chapitre}{La BI traditionnelle reste manuelle, coûteuse et
spécialisée. Les LLM et les agents, couplés au protocole MCP, offrent les briques
d'une automatisation généralisée — à condition d'en maîtriser les faiblesses
(hallucinations, fiabilité). Le chapitre suivant traduit ce constat en une
architecture concrète.}
